\subsection*{Acknowledgments}
The independent referee reports of September 9, 2026 supplied the
quadratic support-overlap benchmark, the cubic physical count-testing
comparison, and the fixed-offset four-window construction. We thank the
referees for these calculations and for requesting that exact-family
inversion, unknown smooth remainders, and joint timing loss be
distinguished. Source identities and the detailed responses are
provided separately from this article. The subsequent smooth-envelope
splice, its sharp minimax comparison and the positive-window equality
were supplied in the AI referee-style memorandum~\cite{A2ReviewV8};
Section~\ref{sec:v9-envelope-minimax} includes their full proof and
interior-slack hypothesis. These memoranda were obtained by the author
for independent review; they are not journal-commissioned reports. The subsequent memorandum
\cite{A2ReviewV9} identified the closest deautoconvolution comparison,
the distinction between reconstruction and compatibility validation,
and the finite-preparation issue quantified by its direct-variance
calculation. The acquisition theorem here addresses that issue with
a separate regularization and sampling argument. The later memorandum
\cite{A2ReviewV10} supplied the exact monomial linearization,
the fixed-class $C^2$ instability benchmark, a mixed-norm estimate,
and the calibration sensitivity of normalized nodal means. We acknowledge
these calculations explicitly. The Abel-flux discrepancy, its integrated
sampling allocation, and the smooth-class calibration procedure are
proved separately here. They do not assert an optimal full-profile
experiment or turn abstract profile alternatives into billiard tables.
The memorandum~\cite{A2ReviewV11} supplied the two-derivative
integrated-flux gain, the higher-order positive-node rate and the
conditional modulus. Section~\ref{sec:v12-regularized-observation}
acknowledges and proves these improvements. Its distinction between
closed correctness objections and the geometric significance question
motivated the separate two-contact inverse and physical realization
in Sections~\ref{sec:v12-two-contact}--\ref{sec:v12-physical-image}.

The memorandum~\cite{A2ReviewV12} extracted and proved the two-flight
independent-contact block and observed that the physical finite-window
design can use $j_0=2$. Section~\ref{sec:v13-two-flight} includes the
complete action, twist and moment proof and its observation consequence.
The finite-order coordinate comparison there makes explicit how that
result relates to the limiting inverse. We acknowledge this source;
these author-requested AI memoranda are not journal editorial decisions.

The September 10, 2026 referee benchmark for this version supplied the
analytic mixed-boundary equation, the physical projection formula and
the stable--unstable interpretation of the relative amplitude. The
September 28 rereview emphasized the finite-family quantifier and the
area-normalization convention. Section~\ref{sec:v14-normal-form}
includes the comparator with its proof and precise scope.
Section~\ref{sec:v14-observability} addresses the observation questions
through a rank characterization and a free-area physical construction.
Both reports and their frozen source identities are retained with the
revision. They are author-requested AI referee-style assessments, not
journal decisions or formal proof certificates.

